
\documentclass{article}
\usepackage{colortbl}
\usepackage{makecell}
\usepackage{multirow}
\usepackage{supertabular}

\begin{document}

\newcounter{utterance}

\centering \large Interaction Transcript for game `dond', experiment `coop\_en', episode 12 with Dialogue Architects\_\_qwen3.5{-}9b\_\_SCoRe\_Qwen3.5{-}9B.
\vspace{24pt}

{ \footnotesize  \setcounter{utterance}{1}
\setlength{\tabcolsep}{0pt}
\begin{supertabular}{c@{$\;$}|p{.15\linewidth}@{}p{.15\linewidth}p{.15\linewidth}p{.15\linewidth}p{.15\linewidth}p{.15\linewidth}}
    \# & $\;$A & \multicolumn{4}{c}{Game Master} & $\;\:$B\\
    \hline

    \theutterance \stepcounter{utterance}  
    & & \multicolumn{4}{p{0.6\linewidth}}{
        \cellcolor[rgb]{0.9,0.9,0.9}{
            \makecell[{{p{\linewidth}}}]{
                \texttt{\tiny{[P1$\langle$GM]}}
                \texttt{You are playing a cooperative negotiation game in which you have to agree on how to divide a set of items among you and another player.} \\
\\ 
\texttt{Rules:} \\
\texttt{(a) You and the other player are given a set of items. Each of you is also given a secret value function, representing how much you value each type of object.} \\
\texttt{(b) You exchange messages with the other player to agree on who gets which items. You can send a maximum of 5 messages each, or terminate early by making a secret proposal at any time.} \\
\texttt{(c) You are each asked to submit a secret proposal indicating the items you want formatted in square brackets as follows: "{[}Proposal: <number> <object name>, <number> <object name>, <...>{]}"} \\
\texttt{(d) If your proposals are complementary, i.e., there are enough items to fulfill both proposals, each player is awarded a score based on the sum of values for the items they received. Otherwise, both of you get zero points.} \\
\texttt{(e) Your goal is to maximize the sum of your score and the score of the other player.} \\
\\ 
\texttt{Let us start.} \\
\\ 
\texttt{The set of available items is:} \\
\\ 
\texttt{1 dice, 3 rings, 2 rulers} \\
\\ 
\texttt{Your secret value function is:} \\
\\ 
\texttt{dice: 0, ring: 2, ruler: 2} \\
\\ 
\texttt{IMPORTANT: Your messages, unless it is the secret proposal, are directly transmitted to the other player, so do not include any response to the rules or text announcing your message. To make a secret proposal, use the indicated format. Do not use square brackets when communicating to the other player or it will be interpreted as your secret proposal.} \\
            }
        }
    }
    & \\ \\

    \theutterance \stepcounter{utterance}  
    & \multicolumn{4}{p{0.6\linewidth}}{
        \cellcolor[rgb]{0.9,0.9,0.9}{
            \makecell[{{p{\linewidth}}}]{
                \texttt{\tiny{[P1$\rangle$GM]}}
                \texttt{I understand the rules and my valuation. I value rings and rulers at 2 points each, and dice at 0. To maximize the total score, I need to ensure the other player gets the dice (which I don't value) and I get the rings and rulers.} \\
\\ 
\texttt{Since I cannot see the other player's values, I will propose a division that gives them the item I value least (the dice) and keeps the high{-}value items for myself, assuming they might value the dice or simply want to avoid conflict. However, a better strategy in a cooperative game where we want to maximize the *sum* is to propose a split that leaves no items unassigned and ensures we don't clash.} \\
\\ 
\texttt{Given the items: 1 dice, 3 rings, 2 rulers.} \\
\texttt{My ideal split: I take 3 rings and 2 rulers (Total value: 10). They take 1 dice.} \\
\texttt{If they value the dice highly, they might reject this if I don't offer them something else, but I have nothing else to offer.} \\
\texttt{Alternatively, I could propose a 50/50 split of the high{-}value items if I think they value them too, but I don't know that.} \\
\texttt{The safest bet to ensure a non{-}zero sum (which is the primary constraint) while maximizing my own score is to take everything I value and give them the item I don't value.} \\
\\ 
\texttt{Proposal:} \\
\texttt{I want the 3 rings and 2 rulers. I will give the dice to the other player.} \\
\\ 
\texttt{Message to other player:} \\
\texttt{I propose you take the 1 dice, and I take the 3 rings and 2 rulers. This way we both get items without conflict.} \\
            }
        }
    }
    & & \\ \\

    \theutterance \stepcounter{utterance}  
    & & & \multicolumn{2}{p{0.3\linewidth}}{
        \cellcolor[rgb]{0.9,0.9,0.9}{
            \makecell[{{p{\linewidth}}}]{
                \texttt{\tiny{[GM$|$GM]}}
                \texttt{continue} \\
            }
        }
    }
    & & \\ \\

    \theutterance \stepcounter{utterance}  
    & & & \multicolumn{2}{p{0.3\linewidth}}{
        \cellcolor[rgb]{0.9,0.9,0.9}{
            \makecell[{{p{\linewidth}}}]{
                \texttt{\tiny{[GM$|$GM]}}
                \texttt{I understand the rules and my valuation. I value rings and rulers at 2 points each, and dice at 0. To maximize the total score, I need to ensure the other player gets the dice (which I don't value) and I get the rings and rulers.} \\
\\ 
\texttt{Since I cannot see the other player's values, I will propose a division that gives them the item I value least (the dice) and keeps the high{-}value items for myself, assuming they might value the dice or simply want to avoid conflict. However, a better strategy in a cooperative game where we want to maximize the *sum* is to propose a split that leaves no items unassigned and ensures we don't clash.} \\
\\ 
\texttt{Given the items: 1 dice, 3 rings, 2 rulers.} \\
\texttt{My ideal split: I take 3 rings and 2 rulers (Total value: 10). They take 1 dice.} \\
\texttt{If they value the dice highly, they might reject this if I don't offer them something else, but I have nothing else to offer.} \\
\texttt{Alternatively, I could propose a 50/50 split of the high{-}value items if I think they value them too, but I don't know that.} \\
\texttt{The safest bet to ensure a non{-}zero sum (which is the primary constraint) while maximizing my own score is to take everything I value and give them the item I don't value.} \\
\\ 
\texttt{Proposal:} \\
\texttt{I want the 3 rings and 2 rulers. I will give the dice to the other player.} \\
\\ 
\texttt{Message to other player:} \\
\texttt{I propose you take the 1 dice, and I take the 3 rings and 2 rulers. This way we both get items without conflict.} \\
            }
        }
    }
    & & \\ \\

    \theutterance \stepcounter{utterance}  
    & & \multicolumn{4}{p{0.6\linewidth}}{
        \cellcolor[rgb]{0.9,0.9,0.9}{
            \makecell[{{p{\linewidth}}}]{
                \texttt{\tiny{[GM$\rangle$P2]}}
                \texttt{You are playing a cooperative negotiation game in which you have to agree on how to divide a set of items among you and another player.} \\
\\ 
\texttt{Rules:} \\
\texttt{(a) You and the other player are given a set of items. Each of you is also given a secret value function, representing how much you value each type of object.} \\
\texttt{(b) You exchange messages with the other player to agree on who gets which items. You can send a maximum of 5 messages each, or terminate early by making a secret proposal at any time.} \\
\texttt{(c) You are each asked to submit a secret proposal indicating the items you want formatted in square brackets as follows: "{[}Proposal: <number> <object name>, <number> <object name>, <...>{]}"} \\
\texttt{(d) If your proposals are complementary, i.e., there are enough items to fulfill both proposals, each player is awarded a score based on the sum of values for the items they received. Otherwise, both of you get zero points.} \\
\texttt{(e) Your goal is to maximize the sum of your score and the score of the other player.} \\
\\ 
\texttt{Let us start.} \\
\\ 
\texttt{The set of available items is:} \\
\\ 
\texttt{1 dice, 3 rings, 2 rulers} \\
\\ 
\texttt{Your secret value function is:} \\
\\ 
\texttt{dice: 3, ring: 1, ruler: 2} \\
\\ 
\texttt{IMPORTANT: Your messages, unless it is the secret proposal, are directly transmitted to the other player, so do not include any response to the rules or text announcing your message. To make a secret proposal, use the indicated format. Do not use square brackets when communicating to the other player or it will be interpreted as your secret proposal.} \\
\\ 
\\ 
\texttt{I understand the rules and my valuation. I value rings and rulers at 2 points each, and dice at 0. To maximize the total score, I need to ensure the other player gets the dice (which I don't value) and I get the rings and rulers.} \\
\\ 
\texttt{Since I cannot see the other player's values, I will propose a division that gives them the item I value least (the dice) and keeps the high{-}value items for myself, assuming they might value the dice or simply want to avoid conflict. However, a better strategy in a cooperative game where we want to maximize the *sum* is to propose a split that leaves no items unassigned and ensures we don't clash.} \\
\\ 
\texttt{Given the items: 1 dice, 3 rings, 2 rulers.} \\
\texttt{My ideal split: I take 3 rings and 2 rulers (Total value: 10). They take 1 dice.} \\
\texttt{If they value the dice highly, they might reject this if I don't offer them something else, but I have nothing else to offer.} \\
\texttt{Alternatively, I could propose a 50/50 split of the high{-}value items if I think they value them too, but I don't know that.} \\
\texttt{The safest bet to ensure a non{-}zero sum (which is the primary constraint) while maximizing my own score is to take everything I value and give them the item I don't value.} \\
\\ 
\texttt{Proposal:} \\
\texttt{I want the 3 rings and 2 rulers. I will give the dice to the other player.} \\
\\ 
\texttt{Message to other player:} \\
\texttt{I propose you take the 1 dice, and I take the 3 rings and 2 rulers. This way we both get items without conflict.} \\
            }
        }
    }
    & \\ \\

    \theutterance \stepcounter{utterance}  
    & & & \multicolumn{4}{p{0.6\linewidth}}{
        \cellcolor[rgb]{0.9,0.9,0.9}{
            \makecell[{{p{\linewidth}}}]{
                \texttt{\tiny{[GM$\langle$P2]}}
                \texttt{I propose you take the 1 dice, and I take the 3 rings and 2 rulers. This way we both get items without conflict.} \\
            }
        }
    }
     \\ \\

    \theutterance \stepcounter{utterance}  
    & & & \multicolumn{2}{p{0.3\linewidth}}{
        \cellcolor[rgb]{0.9,0.9,0.9}{
            \makecell[{{p{\linewidth}}}]{
                \texttt{\tiny{[GM$|$GM]}}
                \texttt{continue} \\
            }
        }
    }
    & & \\ \\

    \theutterance \stepcounter{utterance}  
    & & & \multicolumn{2}{p{0.3\linewidth}}{
        \cellcolor[rgb]{0.9,0.9,0.9}{
            \makecell[{{p{\linewidth}}}]{
                \texttt{\tiny{[GM$|$GM]}}
                \texttt{I propose you take the 1 dice, and I take the 3 rings and 2 rulers. This way we both get items without conflict.} \\
            }
        }
    }
    & & \\ \\

    \theutterance \stepcounter{utterance}  
    & & \multicolumn{4}{p{0.6\linewidth}}{
        \cellcolor[rgb]{0.9,0.9,0.9}{
            \makecell[{{p{\linewidth}}}]{
                \texttt{\tiny{[P1$\langle$GM]}}
                \texttt{I propose you take the 1 dice, and I take the 3 rings and 2 rulers. This way we both get items without conflict.} \\
            }
        }
    }
    & \\ \\

    \theutterance \stepcounter{utterance}  
    & \multicolumn{4}{p{0.6\linewidth}}{
        \cellcolor[rgb]{0.9,0.9,0.9}{
            \makecell[{{p{\linewidth}}}]{
                \texttt{\tiny{[P1$\rangle$GM]}}
                \texttt{That sounds like a fair and efficient split! Since I value the dice at 0, I am happy to give them to you, and I will take the 3 rings and 2 rulers which I value at 2 points each.} \\
\\ 
\texttt{This division ensures we don't conflict over any items, so we both get a non{-}zero score.} \\
\\ 
\texttt{{[}Proposal: 3 ring, 2 ruler{]}} \\
            }
        }
    }
    & & \\ \\

    \theutterance \stepcounter{utterance}  
    & & & \multicolumn{2}{p{0.3\linewidth}}{
        \cellcolor[rgb]{0.9,0.9,0.9}{
            \makecell[{{p{\linewidth}}}]{
                \texttt{\tiny{[GM$|$GM]}}
                \texttt{proposal} \\
            }
        }
    }
    & & \\ \\

    \theutterance \stepcounter{utterance}  
    & & & \multicolumn{2}{p{0.3\linewidth}}{
        \cellcolor[rgb]{0.9,0.9,0.9}{
            \makecell[{{p{\linewidth}}}]{
                \texttt{\tiny{[GM$|$GM]}}
                [0, 3, 2]
            }
        }
    }
    & & \\ \\

    \theutterance \stepcounter{utterance}  
    & & \multicolumn{4}{p{0.6\linewidth}}{
        \cellcolor[rgb]{0.9,0.9,0.9}{
            \makecell[{{p{\linewidth}}}]{
                \texttt{\tiny{[GM$\rangle$P2]}}
                \texttt{The other player has submitted a secret proposal. It is now your turn to submit a proposal of your own.} \\
            }
        }
    }
    & \\ \\

    \theutterance \stepcounter{utterance}  
    & & & \multicolumn{4}{p{0.6\linewidth}}{
        \cellcolor[rgb]{0.9,0.9,0.9}{
            \makecell[{{p{\linewidth}}}]{
                \texttt{\tiny{[GM$\langle$P2]}}
                \texttt{{[}Proposal: 3 ring, 2 ruler{]}} \\
            }
        }
    }
     \\ \\

    \theutterance \stepcounter{utterance}  
    & & & \multicolumn{2}{p{0.3\linewidth}}{
        \cellcolor[rgb]{0.9,0.9,0.9}{
            \makecell[{{p{\linewidth}}}]{
                \texttt{\tiny{[GM$|$GM]}}
                \texttt{proposal} \\
            }
        }
    }
    & & \\ \\

    \theutterance \stepcounter{utterance}  
    & & & \multicolumn{2}{p{0.3\linewidth}}{
        \cellcolor[rgb]{0.9,0.9,0.9}{
            \makecell[{{p{\linewidth}}}]{
                \texttt{\tiny{[GM$|$GM]}}
                [0, 3, 2]
            }
        }
    }
    & & \\ \\

    \theutterance \stepcounter{utterance}  
    & & & \multicolumn{2}{p{0.3\linewidth}}{
        \cellcolor[rgb]{0.9,0.9,0.9}{
            \makecell[{{p{\linewidth}}}]{
                \texttt{\tiny{[GM$|$GM]}}
                [[0, 3, 2], [0, 3, 2]]
            }
        }
    }
    & & \\ \\

\end{supertabular}
}

\end{document}
